\appendix
\section{Tái lập và đối chiếu hiện thực}

Từ thư mục gốc của dự án, thực hiện các lệnh sau. Cần Python 3.13 và uv;
phiên bản phụ thuộc được khóa trong \code{uv.lock}.
\begin{lstlisting}[numbers=none,basicstyle=\ttfamily\footnotesize]
uv sync --frozen
uv run main.py
uv run python -m unittest discover -s pipes_problem -t .
uv run python -m pipes_problem.benchmark
\end{lstlisting}

Thí nghiệm mới tạo thư mục dưới \code{results/}, có tên bắt đầu bằng
\code{benchmark-} và mốc thời gian. Mỗi thư mục gồm \code{benchmark.csv},
\code{manifest.json} và \code{summary.md}. Các giá trị
thời gian mới có thể thay đổi theo máy. Bảng trong báo cáo này lấy từ bằng
chứng đã lưu ngày 22/10/2026, không tự động thay bằng kết quả mới.

\begin{table}[H]\centering\small
\begin{tabularx}{\textwidth}{@{}lX@{}}\toprule
Cột CSV & Ý nghĩa \\\midrule
\code{puzzle_id} & Tên bản đồ không có đuôi .json.\\
\code{grid_size} & Số ô trên một cạnh.\\
\code{algorithm} & \code{dfs} hoặc \code{greedy_best_first}.\\
\code{seed} & Seed của thuật toán ngẫu nhiên; trống với hai thuật toán hiện tại.\\
\code{status} & solved, unsolvable, limit\_reached, cancelled hoặc error.\\
\code{time_ms} & Thời gian tính toán được đo theo mili giây.\\
\code{peak_memory_kib} & Đỉnh cấp phát Python, để trống khi không đo được.\\\bottomrule
\end{tabularx}
\caption{Bảy cột lưu tự động; tham số và số mở rộng nằm trong manifest.}
\end{table}

\begin{table}[H]\centering\small
\begin{tabularx}{\textwidth}{@{}>{\raggedright\arraybackslash}p{.41\textwidth}>{\raggedright\arraybackslash}X@{}}\toprule
Thành phần & Vị trí trong mã nguồn \\\midrule
Trạng thái, sinh hướng và cắt tỉa & \code{Search}, \code{orientations()}, \code{compatible()} trong \code{search.py}.\\
DFS và GBFS & \code{Search.dfs()} và \code{Search.greedy()}.\\
Kiểm tra đích & \code{inspect_board()} trong \code{model.py}, \code{validate_solution()} trong \code{search.py}.\\
Thời gian và bộ nhớ & \code{solve()}, \code{measure()} và \code{worker.py}.\\
Giao diện và sự kiện & \code{advance_search()}, \code{draw_board()} và \code{draw_sidebar()} trong \code{app.py}.\\
Thực nghiệm lô & \code{benchmark.py} và \code{benchmark.json}.\\\bottomrule
\end{tabularx}
\caption{Đối chiếu nội dung báo cáo với hiện thực.}
\end{table}

Mỗi manifest lưu chính xác bản đồ, hash, môi trường, cấu hình, số lần lặp và
số mở rộng. Hash nguồn của lô dùng trong báo cáo là:\par
{\footnotesize\ttfamily\nolinkurl{4a9c58c27c2c8a0ad3d9e30ce02742b057c23d2c0fcdefef6774a6fa9ca8d93f}}.
